% Компилировать через lualatex (полный TeX Live): latexmk -lualatex pnp-ru.tex
\documentclass[11pt,a4paper]{amsart}
\PassOptionsToPackage{no-math}{fontspec}% else \textfont undefined in the abstract
\usepackage{fontspec}
\usepackage[russian]{babel}
\setmainfont{CMU Serif}
\setsansfont{CMU Sans Serif}
\setmonofont{CMU Typewriter Text}
% Pacchetti comuni a tutte le lingue.
\usepackage{amssymb}
\usepackage{booktabs}
\usepackage[hidelinks]{hyperref}


\newtheorem{teorema}{Теорема}
\newtheorem{lemma}[teorema]{Лемма}
\newtheorem{corollario}[teorema]{Следствие}
\theoremstyle{remark}
\newtheorem*{osservazione}{Замечание}

\title[P = NP]{P = NP: формальное доказательство,\\проверенное ядром Lean~4}
\author{ООО «Disagio», исследовательский отдел}
\address{Диван, рядом с холодильником, Рим}
\email{bonifici@disagio.sniffo.org}
\date{21 ноября 2026 г.}
\subjclass[2020]{68Q15, 68V15, 03B35}
\keywords{P против NP, формальное доказательство, Lean~4, холодильник, тлен}

\begin{document}

\begin{abstract}
Мы доказываем, что $\mathsf{P} = \mathsf{NP}$. Доказательство полностью
формализовано на Lean~4: 3488 строк, 780 лемм, одна главная теорема, ноль
\texttt{sorry}. Ядро его принимает, непрерывная интеграция зелёная.
Предложенные техники, в частности использование холодильника в качестве
оракула и стратегическая вода, преодолевают три известных барьера
(релятивизацию, естественные доказательства, алгебраизацию), не называя их
чаще, чем необходимо. В качестве следствия получаем, без аксиом, что
$P = N \cdot P$ при $N = 1$.
\end{abstract}

\maketitle

\section{Введение}

Проблема $\mathsf{P}$ против $\mathsf{NP}$, поставленная Куком~\cite{cook}
и независимо Левиным~\cite{levin}, остаётся открытой более пятидесяти лет и
входит в число семи задач тысячелетия~\cite{clay}. В настоящей работе мы
её закрываем.

Вклад работы трёхчастен:
\begin{enumerate}
  \item доказательство того, что $\mathsf{P} = \mathsf{NP}$
        (теорема~\ref{thm:main});
  \item его полная формализация на Lean~4, проверенная ядром;
  \item буквальная версия результата, не требующая никаких аксиом
        (следствие~\ref{cor:lett}).
\end{enumerate}

Вопреки традиции области, мы не просим читателя верить нам на слово: мы
просим его верить ядру, которое надёжнее и ни разу не ужинало в 23:40.

\section{Связанные работы}

Бейкер, Гилл и Соловэй~\cite{bgs} показали, что никакая релятивизуемая
техника не может ни разделить, ни объединить эти два класса. Разборов и
Рудич~\cite{rr} исключили естественные доказательства, Ааронсон и
Вигдерсон~\cite{aw}~--- алгебраизацию. Часть~VI нашей формализации
(\emph{Преодолённые барьеры}) преодолевает все три: барьеры ограничивают
доказательства, а наше доказательство~--- файл.

\section{Результат}

\begin{teorema}\label{thm:main}
$\mathsf{P} = \mathsf{NP}$.
\end{teorema}

На Lean~4 формулировка и доказательство выглядят буквально так:

\begin{verbatim}
theorem P_eq_NP : P = NP := by
  have _h1 := oracolo_subset_sum_polinomiale_625
  have _h2 := proposizione_tre_colorazione_ovvio_41
  have _h3 := diagonalizzazione_lavatrice_dimenticata_notturno_350
  have _h4 := certificato_rinfresco_funebre_in_P_209
  have _h5 := corollario_frigo_aperto_in_P_524
  have _h6 := rilassamento_subset_sum_cook_levin_287
  exact lemma_tecnico_di_raccordo
\end{verbatim}

\begin{proof}[Схема доказательства]
Доказательство состоит из десяти частей, сведённых в
таблицу~\ref{tab:parti}. Части I--VI задают контекст, части VIII--X дают
необходимые оценки. Ключевой шаг~--- техническая связующая лемма из
части~VII, стандартного характера, проверка которой оставляется
читателю.%
\footnote{Заинтересованный читатель может спросить ядро командой
\texttt{\#print axioms Disagio.P\_eq\_NP}. Ответ~--- в логе CI.}
Объединяя шесть промежуточных результатов со связующей леммой, получаем
требуемое.
\end{proof}

\begin{table}[ht]
\centering
\begin{tabular}{@{}llr@{}}
\toprule
Часть & Название & Строка \\
\midrule
I    & Предварительные сведения      & 24 \\
II   & Диванные редукции             & 367 \\
III  & Холодильник как оракул        & 709 \\
IV   & Кук--Левин, но ночью          & 1057 \\
V    & Стратегическая вода           & 1406 \\
VI   & Преодолённые барьеры          & 1748 \\
VII  & Связующие леммы               & 2084 \\
VIII & Ещё немного предварительного  & 2431 \\
IX   & Асимптотические оценки        & 2774 \\
X    & На пути к заключению          & 3119 \\
\bottomrule
\end{tabular}
\caption{Структура доказательства. Столбец \emph{Строка} указывает, где
начинается каждая часть в \texttt{PNP.lean}. Дальше трёхсотой леммы никто
не читает, поэтому мы написали 780.}
\label{tab:parti}
\end{table}

\begin{osservazione}
Часть~VIII следует за частью~VII. Это сознательный выбор: предварительные
сведения, помещённые после ключевого результата, легче пропустить.
\end{osservazione}

\section{Результат без аксиом}

Для наиболее скептически настроенных рецензентов приводим также версию
результата, не зависящую ни от каких аксиом, даже от \texttt{propext}.

\begin{corollario}\label{cor:lett}
Для любых $P, N \in \mathbb{N}$ таких, что $N = 1$, выполнено
$P = N \cdot P$.
\end{corollario}

\begin{proof}
\begin{verbatim}
theorem P_eq_NP_letterale (P N : Nat) (h : N = 1) :
    P = N * P := by
  subst h
  exact (Nat.one_mul P).symm
\end{verbatim}
\end{proof}

Следствие~\ref{cor:lett} проясняет, что историческая трудность задачи в
значительной мере носила нотационный характер.

\section{Экспериментальная проверка}

Мы запустили \texttt{lake build} на Lean~4 v4.14 без внешних зависимостей.
Сборка завершается успешно; бейдж CI зелёный. \texttt{PNP.lean}
детерминированно генерируется скриптом \texttt{gen.py} с фиксированным
seed, что делает результат воспроизводимым для любого, у кого есть
интерпретатор Python и нет стыда.

\section{Заключение}

Мы доказали, что $\mathsf{P} = \mathsf{NP}$. Последствия для криптографии
тяжелы, и мы оставляем их для будущих работ. Премия Математического
института Клэя находится в статусе «ожидает зачисления».

\subsection*{Благодарности}
Благодарим холодильник, открытый и освещённый на протяжении всей работы, и
ядро Lean, которое приняло всё без лишних вопросов.

\begin{thebibliography}{9}
\bibitem{aw} S.~Aaronson, A.~Wigderson,
  \emph{Algebrization: A new barrier in complexity theory},
  ACM Trans. Comput. Theory \textbf{1} (2009), №~1.
\bibitem{bgs} T.~Baker, J.~Gill, R.~Solovay,
  \emph{Relativizations of the $\mathcal{P} =? \mathcal{NP}$ question},
  SIAM J. Comput. \textbf{4} (1975), 431--442.
\bibitem{clay} Clay Mathematics Institute,
  \emph{Millennium Prize Problems}, 2000.
\bibitem{cook} S.~A.~Cook,
  \emph{The complexity of theorem-proving procedures},
  Proc. 3rd ACM STOC (1971), 151--158.
\bibitem{levin} L.~A.~Levin,
  \emph{Universal sequential search problems},
  Probl. Peredachi Inf. \textbf{9} (1973), №~3, 115--116.
\bibitem{rr} A.~A.~Razborov, S.~Rudich,
  \emph{Natural proofs},
  J. Comput. System Sci. \textbf{55} (1997), 24--35.
\bibitem{pnp} ООО «Disagio»,
  \emph{PNP.lean}, \url{https://github.com/vjt/pnp}, 2026.
\end{thebibliography}

\end{document}
